\chapter{研究内容与方法}
\section{研究方案}
图应先在正文提及，再紧随相关文字放置。图~\ref{fig:example} 是图题与分章编号示例。

\begin{figure}[htbp]
  \centering
  \fbox{\parbox[c][5cm][c]{6.67cm}{\centering 请在此放入清晰的图片}}
  \caption{示例图题}\label{fig:example}
\end{figure}

表~\ref{tab:example} 为三线表示例。表题在表格上方，表与正文同宽，数字按列对齐。
\begin{table}[htbp]
  \centering
  \caption{示例数据表}\label{tab:example}
  \zihao{5}
  \begin{tabular*}{\textwidth}{@{\extracolsep{\fill}}lcc@{}}
    \toprule[1.5pt]
    项目 & 数值 / 单位 & 说明 \\
    \midrule[1pt]
    样本 A & 10 & 示例数据 \\
    样本 B & 12 & 示例数据 \\
    \bottomrule[1.5pt]
  \end{tabular*}
\end{table}

公式单独成行，按章编号为括号中的“章号-序号”。
\begin{equation}\label{eq:example}
  y=ax+b
\end{equation}
\noindent 式中，$y$ 为响应量；$x$ 为自变量；$a$、$b$ 为参数。
